%=============================================================================!
% 3.9  WORK ALONG A PATH IN SPACE                                             !
%=============================================================================!
\section{Work along a path in space}
\label{sec:en-paths}

On a line, every force that depends only on position has a potential
energy, because the integral of \cref{eq:en-potential} can always be
done. In a plane or in space a body can travel from one point to another
by many routes, and the work may differ from route to route. This section
defines the work along a curved path, the tool we need, and shows that
the route can matter.

\subsection*{The line integral}

\Cref{sec:en-track} found the work of the weight along a curved wire by
cutting the wire into short straight chords. The same construction works
for any force, constant or not, and the box below turns it into a
definition.

\begin{toolbox}[Line integrals]
Let a path $C$ run from a point $A$ to a point $B$, and let a force
$\vb{F}(\vb{r})$ depend on position. Mark points $\vb{r}_0 = A$,
$\vb{r}_1, \dots, \vb{r}_n = B$ in order along the path, and join each to
the next by a straight chord, $\Delta\vb{r}_k = \vb{r}_k -
\vb{r}_{k-1}$. On a chord short enough that the force hardly changes
along it, the force does about $\vb{F}(\bar{\vb{r}}_k)\cdot\Delta
\vb{r}_k$ of work by \cref{eq:en-work-constant}, where $\bar{\vb{r}}_k =
\frac12(\vb{r}_{k-1} + \vb{r}_k)$ is the midpoint of the chord. The work
along the path is the limit of the sum as the chords shrink,
\begin{equation*}
   W[C] = \lim\sum_{k=1}^{n}\vb{F}(\bar{\vb{r}}_k)\cdot\Delta\vb{r}_k
   = \int_C \vb{F}\cdot\dd\vb{r} ,
\end{equation*}
called the line integral of $\vb{F}$ along $C$; the square brackets in
$W[C]$ record that the work belongs to the whole path $C$.
Three properties follow from the sum. Reversing the direction of travel
reverses every chord, $\Delta\vb{r}_k \to -\Delta\vb{r}_k$, and so changes
the sign of the work. The sum uses only the points of the path and their
order, so it does not matter how fast the path is traversed. And along a
path made of one path followed by another, the work is the sum of the
works along the two, since the sum over chords splits into two such sums.
For a closed path, one that returns to its start, the integral is written
$\oint_C$.

To compute a line integral, describe the path by a parameter $s$ that
runs from $s_A$ to $s_B$, so that $\vb{r}(s)$ traces the path, and cut
the range of $s$ into short steps $\Delta s$, with $s_k$ the middle of the
$k$th. Each step moves the point by $\Delta\vb{r}_k \approx (\dd\vb{r}/\dd
s)\,\Delta s$, where $\dd\vb{r}/\dd s$, the derivative of each component
of $\vb{r}$, is taken at $s_k$. Then
\begin{align*}
   \int_C \vb{F}\cdot\dd\vb{r}
   &= \lim\sum_k\vb{F}\bigl(\vb{r}(s_k)\bigr)\cdot
      \frac{\dd\vb{r}}{\dd s}\,\Delta s
      && \text{putting in the chords}\\
   &= \int_{s_A}^{s_B}\vb{F}\bigl(\vb{r}(s)\bigr)\cdot
      \frac{\dd\vb{r}}{\dd s}\,\dd s
      && \text{a sum of values times $\Delta s$, \cref{sec:mo-integrating},}
\end{align*}
an ordinary integral in one variable.
For example, take the force $\vb{F} = c\,(-y, x)$ in the plane, with $c$
a constant, and the circle of radius $R$ about the origin, traversed once
anticlockwise: $\vb{r}(s) = R\,(\cos s, \sin s)$ for $s$ from 0 to $2\pi$.
Then
\begin{align*}
   \frac{\dd\vb{r}}{\dd s} &= R\,(-\sin s, \cos s)
      && \text{differentiating each component}\\
   \vb{F}\bigl(\vb{r}(s)\bigr) &= c\,(-R\sin s, R\cos s)
      && \text{setting $x = R\cos s$, $y = R\sin s$}\\
   \vb{F}\cdot\frac{\dd\vb{r}}{\dd s}
   &= cR^2\sin^2 s + cR^2\cos^2 s
      && \text{multiplying components and adding}\\
   &= cR^2
      && \text{since $\sin^2 s + \cos^2 s = 1$,}
\end{align*}
and the line integral is $\int_0^{2\pi} cR^2\,\dd s = 2\pi cR^2$.
\end{toolbox}

Along the motion itself the time $t$ is a natural parameter, and
$\dd\vb{r}/\dd t$ is the velocity. The work-energy theorem
\cref{eq:en-work-energy} therefore reads, in space,
\begin{equation*}
   K_B - K_A = \int_{t_A}^{t_B}\vb{F}\cdot\vb{v}\,\dd t
   = \int_{\text{path}}\vb{F}\cdot\dd\vb{r} ,
\end{equation*}
the line integral of the net force along the path the body actually
follows.

\subsection*{When the route matters}

The weight is the simplest case. With $z$ vertical, $\vb{F} = (0, 0,
-mg)$, each chord contributes $-mg\,\Delta z_k$, and the sum is $-mg\,(z_B
- z_A)$ for every path from $A$ to $B$, as on the wire of
\cref{sec:en-track}.

Friction behaves differently. A box of \SI{5}{\kilogram} is pushed slowly
across a floor with $\mu_{\mathrm{k}} = 0.3$ from $A$ to a point $B$
\SI{5}{\metre} away. Friction always points against the motion with size
$\mu_{\mathrm{k}}mg = \SI{14.71}{\newton}$, so on every chord it does
$-14.71$ newtons times the chord's length, and its work is $-14.71$
newtons times the length of the path. Straight from $A$ to $B$ that is
$-14.71\times5 = \SI{-73.55}{\joule}$; along the other two sides of a
right-angled triangle, of \SI{3}{\metre} and \SI{4}{\metre}, it is
$-14.71\times7 = \SI{-102.97}{\joule}$. Friction depends on the direction of motion, so
this is no great surprise.

More surprising is that a force depending only on position can also do
different work on different routes. The force of the toolbox,
\begin{equation}
   \vb{F}_{\mathrm{sw}}(\vb{r}) = c\,(-y, x) ,
   \label{eq:en-swirl}
\end{equation}
which we call a swirl, has size $c\lvert\vb{r}\rvert$ and points around
the origin, perpendicular to $\vb{r}$, since $(-y, x)\cdot(x, y) = 0$.
Once round a circle it does the work $2\pi cR^2$. A closed path starts and
ends at the same point, so if the work did not depend on the route, the
circle would give the same work as staying put, which is zero. The work of
the swirl therefore depends on the route, even though the force depends
only on where the body is. No force between atoms has this form; we use the
swirl as the plainest example of such a force, and in \cref{sec:en-drift}
to stand for the errors of a faulty force calculation.

\subsection*{The work in code}

On a computer a path is a list of points, and the line integral is the
sum of the toolbox with the force taken at the midpoint of each chord.

\begin{code}
def work(force: Callable[[NDArray], NDArray], path: ArrayLike) -> float:
    path = np.asarray(path, dtype=float)
    chords = np.diff(path, axis=0)  # the step from each point to the next
    midpoints = 0.5 * (path[1:] + path[:-1])
    # F · Δr on every chord, added over the components and the chords
    return float(np.sum(force(midpoints) * chords))
\end{code}

The first argument, \texttt{force}, is itself a function
(\texttt{Callable}): given an array of positions, one row per point, it
returns the forces there, one row each. The path is an array with one row
per point and one column per component. \texttt{np.diff} subtracts each row from the next, giving the
chords $\Delta\vb{r}_k$, one row each, and \texttt{path[1:]} and
\texttt{path[:-1]}, the points without the first and without the last,
average to the midpoints. The force is evaluated at all the midpoints at
once; multiplying by the chords, component by component, and adding
everything gives $\sum_k\vb{F}(\bar{\vb{r}}_k)\cdot\Delta\vb{r}_k$, since
adding the products of the components of one row is its scalar product.
The tests of \texttt{mdlab} confirm that halving the length of every
chord divides the error of this sum by four.

\notebook{03}{bend a path and compare the work of the weight, of friction
and of the swirl}

\begin{selfcheck}
How much work does the swirl of \cref{eq:en-swirl} do on a body that moves
from the origin straight out to the point $(3, 4)$?
\end{selfcheck}
